%% sections/07-verification.tex
%% Verification & Results
%% IEC/IEEE 82079-1:2019 — Instructional Information (§5)
%% ─────────────────────────────────────────────────────────

\section{Verification and Results}
\label{sec:verification}

%%
%% WRITING GUIDE — VERIFICATION:
%%
%%   This section answers: "How does the operator know it worked?"
%%
%%   Include:
%%   - Acceptance criteria table (quantitative pass/fail thresholds)
%%   - Expected sensor plots (screenshot or generated Python plot)
%%   - Video link to recorded test run
%%   - What to do if any criterion fails (pointer to troubleshooting)
%%

\subsection{Acceptance Criteria}
\label{sec:verif:criteria}

The force control setup is considered SUCCESSFUL when ALL of the following
criteria are met simultaneously:

\begin{center}
\begin{tabularx}{\textwidth}{@{} X l l l @{}}
  \toprule
  \rowcolor{PBATableHeader}
  \textcolor{white}{\bfseries Parameter} &
  \textcolor{white}{\bfseries Criterion} &
  \textcolor{white}{\bfseries Measured value} &
  \textcolor{white}{\bfseries Pass / Fail} \\
  \midrule
  Force setpoint tracking error (steady state) &
    $< \pm$\SI{0.3}{\newton} & \rule{1.5cm}{0.4pt} & $\square$ \\
  \rowcolor{PBALightGray}
  Force noise (RMS, no contact) &
    $< \SI{0.1}{\newton}$ & \rule{1.5cm}{0.4pt} & $\square$ \\
  Step response settling time (\SI{5}{\newton} step) &
    $< \SI{150}{\milli\second}$ & \rule{1.5cm}{0.4pt} & $\square$ \\
  \rowcolor{PBALightGray}
  Overshoot (\SI{5}{\newton} step) &
    $< 20\%$ & \rule{1.5cm}{0.4pt} & $\square$ \\
  No axis fault flags during 60\,s continuous run &
    True & \rule{1.5cm}{0.4pt} & $\square$ \\
  \bottomrule
\end{tabularx}
\end{center}

\begin{notebox}
  Fill in the ``Measured value'' and ``Pass / Fail'' columns during
  commissioning and retain a signed copy as the acceptance record.
  Attach the force log CSV file to this document in Appendix~\ref{app:data}.
\end{notebox}


\subsection{Expected Force Response Plot}
\label{sec:verif:plot}

\Cref{fig:force-response} shows the expected force response during a
\SI{5}{\newton} step command with the baseline gains from
\cref{tab:pid-defaults}.

\begin{figure}[H]
  \centering
  \begin{subfigure}[b]{0.48\textwidth}
    %% Replace with: \includegraphics[width=\textwidth]{figures/force-step-response.png}
    \fbox{\parbox{\textwidth}{\centering\vspace{28mm}
      \textcolor{PBAGray}{[PLOT: Force step response — time vs F(N)]}\vspace{28mm}}}
    \caption{Step response: \SI{0}{\newton}~$\to$~\SI{5}{\newton} command.
             Settling time $<$ \SI{150}{\milli\second}.}
    \label{fig:step-response}
  \end{subfigure}
  \hfill
  \begin{subfigure}[b]{0.48\textwidth}
    %% Replace with: \includegraphics[width=\textwidth]{figures/force-noise-floor.png}
    \fbox{\parbox{\textwidth}{\centering\vspace{28mm}
      \textcolor{PBAGray}{[PLOT: Force noise — free-hanging, no contact]}\vspace{28mm}}}
    \caption{Force noise floor during free-hanging Z-axis
             (no contact).  RMS $<$ \SI{0.1}{\newton}.}
    \label{fig:noise-floor}
  \end{subfigure}
  \caption{Expected force control verification plots.
           Generate these from the CSV log using
           \filepath{code/plot\_force\_results.py}.}
  \label{fig:force-response}
\end{figure}

\fignote{Plots generated with Python matplotlib from \filepath{force\_log.csv}.
         Actual results may differ by $\pm$15\% depending on workpiece
         compliance and cable routing.}


\subsection{Video Reference}
\label{sec:verif:video}

The video below shows a complete verification run on the PBA GS-500
demonstrator machine at PBA Systems:

\videolink{https://your-video-url-here.com}{Verification run — force step response and ACS MMI display (2\,min)}


\subsection{Interpreting Results}
\label{sec:verif:interpret}

\begin{itemize}[noitemsep]
  \item \textbf{Excessive overshoot} ($>$20\,\%): Reduce \acsparam{KP} by
        20\,\% increments until overshoot is acceptable.
  \item \textbf{Slow settling} ($>$\SI{200}{\milli\second}): Increase
        \acsparam{KI} by 0.01 steps.
  \item \textbf{High noise floor}: Check grounding (single-point shield),
        sensor cable routing, and AIN differential mode setting in the ACS MMI.
  \item \textbf{Force drifts over time}: Re-zero the sensor (\cref{sec:procedure},
        Step~2) and check the ambient temperature — the load cell has a
        thermal zero drift coefficient of $\approx$\SI{0.3}{\newton\per\kelvin}.
\end{itemize}

If any acceptance criterion fails, refer to \cref{sec:troubleshooting}
before signing off the acceptance record.
